% Options for packages loaded elsewhere
\PassOptionsToPackage{unicode}{hyperref}
\PassOptionsToPackage{hyphens}{url}
%
\documentclass[
  english,
  man]{apa6}
\usepackage{lmodern}
\usepackage{amsmath}
\usepackage{ifxetex,ifluatex}
\ifnum 0\ifxetex 1\fi\ifluatex 1\fi=0 % if pdftex
  \usepackage[T1]{fontenc}
  \usepackage[utf8]{inputenc}
  \usepackage{textcomp} % provide euro and other symbols
  \usepackage{amssymb}
\else % if luatex or xetex
  \usepackage{unicode-math}
  \defaultfontfeatures{Scale=MatchLowercase}
  \defaultfontfeatures[\rmfamily]{Ligatures=TeX,Scale=1}
\fi
% Use upquote if available, for straight quotes in verbatim environments
\IfFileExists{upquote.sty}{\usepackage{upquote}}{}
\IfFileExists{microtype.sty}{% use microtype if available
  \usepackage[]{microtype}
  \UseMicrotypeSet[protrusion]{basicmath} % disable protrusion for tt fonts
}{}
\makeatletter
\@ifundefined{KOMAClassName}{% if non-KOMA class
  \IfFileExists{parskip.sty}{%
    \usepackage{parskip}
  }{% else
    \setlength{\parindent}{0pt}
    \setlength{\parskip}{6pt plus 2pt minus 1pt}}
}{% if KOMA class
  \KOMAoptions{parskip=half}}
\makeatother
\usepackage{xcolor}
\IfFileExists{xurl.sty}{\usepackage{xurl}}{} % add URL line breaks if available
\IfFileExists{bookmark.sty}{\usepackage{bookmark}}{\usepackage{hyperref}}
\hypersetup{
  pdftitle={The Use of Explicit Health Benefits Packages Increases Support for Universal Health Care for People with High Objective Numeracy},
  pdfauthor={Sean Xiaohao Duan1},
  pdflang={en-EN},
  pdfkeywords={Bayesian, statistical modeling, frequentist, applied statistics},
  hidelinks,
  pdfcreator={LaTeX via pandoc}}
\urlstyle{same} % disable monospaced font for URLs
\usepackage{graphicx}
\makeatletter
\def\maxwidth{\ifdim\Gin@nat@width>\linewidth\linewidth\else\Gin@nat@width\fi}
\def\maxheight{\ifdim\Gin@nat@height>\textheight\textheight\else\Gin@nat@height\fi}
\makeatother
% Scale images if necessary, so that they will not overflow the page
% margins by default, and it is still possible to overwrite the defaults
% using explicit options in \includegraphics[width, height, ...]{}
\setkeys{Gin}{width=\maxwidth,height=\maxheight,keepaspectratio}
% Set default figure placement to htbp
\makeatletter
\def\fps@figure{htbp}
\makeatother
\setlength{\emergencystretch}{3em} % prevent overfull lines
\providecommand{\tightlist}{%
  \setlength{\itemsep}{0pt}\setlength{\parskip}{0pt}}
\setcounter{secnumdepth}{-\maxdimen} % remove section numbering
% Make \paragraph and \subparagraph free-standing
\ifx\paragraph\undefined\else
  \let\oldparagraph\paragraph
  \renewcommand{\paragraph}[1]{\oldparagraph{#1}\mbox{}}
\fi
\ifx\subparagraph\undefined\else
  \let\oldsubparagraph\subparagraph
  \renewcommand{\subparagraph}[1]{\oldsubparagraph{#1}\mbox{}}
\fi
% Manuscript styling
\usepackage{upgreek}
\captionsetup{font=singlespacing,justification=justified}

% Table formatting
\usepackage{longtable}
\usepackage{lscape}
% \usepackage[counterclockwise]{rotating}   % Landscape page setup for large tables
\usepackage{multirow}		% Table styling
\usepackage{tabularx}		% Control Column width
\usepackage[flushleft]{threeparttable}	% Allows for three part tables with a specified notes section
\usepackage{threeparttablex}            % Lets threeparttable work with longtable

% Create new environments so endfloat can handle them
% \newenvironment{ltable}
%   {\begin{landscape}\begin{center}\begin{threeparttable}}
%   {\end{threeparttable}\end{center}\end{landscape}}
\newenvironment{lltable}{\begin{landscape}\begin{center}\begin{ThreePartTable}}{\end{ThreePartTable}\end{center}\end{landscape}}

% Enables adjusting longtable caption width to table width
% Solution found at http://golatex.de/longtable-mit-caption-so-breit-wie-die-tabelle-t15767.html
\makeatletter
\newcommand\LastLTentrywidth{1em}
\newlength\longtablewidth
\setlength{\longtablewidth}{1in}
\newcommand{\getlongtablewidth}{\begingroup \ifcsname LT@\roman{LT@tables}\endcsname \global\longtablewidth=0pt \renewcommand{\LT@entry}[2]{\global\advance\longtablewidth by ##2\relax\gdef\LastLTentrywidth{##2}}\@nameuse{LT@\roman{LT@tables}} \fi \endgroup}

% \setlength{\parindent}{0.5in}
% \setlength{\parskip}{0pt plus 0pt minus 0pt}

% \usepackage{etoolbox}
\makeatletter
\patchcmd{\HyOrg@maketitle}
  {\section{\normalfont\normalsize\abstractname}}
  {\section*{\normalfont\normalsize\abstractname}}
  {}{\typeout{Failed to patch abstract.}}
\patchcmd{\HyOrg@maketitle}
  {\section{\protect\normalfont{\@title}}}
  {\section*{\protect\normalfont{\@title}}}
  {}{\typeout{Failed to patch title.}}
\makeatother
\shorttitle{Effect of HBPs on UHC}
\keywords{Bayesian, statistical modeling, frequentist, applied statistics}
\DeclareDelayedFloatFlavor{ThreePartTable}{table}
\DeclareDelayedFloatFlavor{lltable}{table}
\DeclareDelayedFloatFlavor*{longtable}{table}
\makeatletter
\renewcommand{\efloat@iwrite}[1]{\immediate\expandafter\protected@write\csname efloat@post#1\endcsname{}}
\makeatother
\usepackage{lineno}

\linenumbers
\usepackage{csquotes}
\ifxetex
  % Load polyglossia as late as possible: uses bidi with RTL langages (e.g. Hebrew, Arabic)
  \usepackage{polyglossia}
  \setmainlanguage[]{english}
\else
  \usepackage[shorthands=off,main=english]{babel}
\fi
\ifluatex
  \usepackage{selnolig}  % disable illegal ligatures
\fi

\title{The Use of Explicit Health Benefits Packages Increases Support for Universal Health Care for People with High Objective Numeracy}
\author{Sean Xiaohao Duan\textsuperscript{1}}
\date{}


\authornote{

This paper was prepared as part of a semester-long, PhD-level seminar at the University of Missouri on Bayesian methods in psychology.

}

\affiliation{\vspace{0.5cm}\textsuperscript{1} University of Missouri}

\abstract{
Bayesian methods are becoming increasingly used in applied psychology research. Previous researchers have thoroughly written about much of the details already, including the philosophy underlying Bayesian methods, computational issues associated with Bayesian model estimation, Bayesian model development and summary, and the role of Bayesian methods in the so-called replication crisis. We do not attempt to rehash these topics in this paper, but instead seek to provide case studies comparing the use of frequentist methods to the use of Bayesian methods in applied psychological research. These case studies are intended to ``illustrate by example'' the ways that Bayesian modeling differs from frequentist modeling, and the differing conclusions that one may arrive at using the two methods. The intended audience is applied psychology researchers who have been trained in the traditional frequentist framework, and who are curious about how statistical results might look different in a Bayesian context. Along with our case studies, we provide general opinions and guidance on the use of Bayesian methods in applied psychology research.
}



\begin{document}
\maketitle

\hypertarget{introduction}{%
\section{Introduction}\label{introduction}}

\hypertarget{method}{%
\section{Method}\label{method}}

\hypertarget{participants}{%
\subsection{Participants}\label{participants}}

Our participants were 189 students enrolled in a Psychology course at a Midwestern University. Our participants were primarily white (76\%), female (68\%), and freshmen (80\%); further demographic information can be found in {[}Table here{]}. Participants received course credit for participation in the study.

\begin{tabular}{l|l|l|l|l}
\hline
**Characteristic** & **Active Intervention**, N = 60 & **No Intervention**, N = 62 & **Passive Intervention**, N = 63 & **p-value**\\
\hline
\_\_Age\_\_ &  &  &  & 0.10\\
\hline
18 & 34 / (57\%) & 46 / (74\%) & 44 / (70\%) & \\
\hline
19 & 19 / (32\%) & 10 / (16\%) & 16 / (25\%) & \\
\hline
20 & 7 / (12\%) & 3 / (4.8\%) & 2 / (3.2\%) & \\
\hline
21 & 0 / (0\%) & 2 / (3.2\%) & 1 / (1.6\%) & \\
\hline
23 & 0 / (0\%) & 1 / (1.6\%) & 0 / (0\%) & \\
\hline
\_\_Sex\_\_ &  &  &  & 0.7\\
\hline
Female & 36 / (60\%) & 41 / (66\%) & 42 / (67\%) & \\
\hline
Male & 24 / (40\%) & 21 / (34\%) & 21 / (33\%) & \\
\hline
\_\_Race\_\_ &  &  &  & >0.9\\
\hline
Caucasian/White & 46 / (77\%) & 49 / (79\%) & 48 / (76\%) & \\
\hline
African-American/Black & 5 / (8.3\%) & 5 / (8.1\%) & 6 / (9.5\%) & \\
\hline
Asian/Pacific Islander & 3 / (5.0\%) & 3 / (4.8\%) & 3 / (4.8\%) & \\
\hline
Caucasian/White,Asian/Pacific Islander & 2 / (3.3\%) & 4 / (6.5\%) & 1 / (1.6\%) & \\
\hline
Caucasian/White,African-American/Black & 1 / (1.7\%) & 1 / (1.6\%) & 1 / (1.6\%) & \\
\hline
 & 2 / (3.3\%) & 0 / (0\%) & 0 / (0\%) & \\
\hline
Other & 1 / (1.7\%) & 0 / (0\%) & 1 / (1.6\%) & \\
\hline
African-American/Black,American Indian/Alaska Native & 0 / (0\%) & 0 / (0\%) & 1 / (1.6\%) & \\
\hline
African-American/Black,Hispanic/Latino(a) & 0 / (0\%) & 0 / (0\%) & 1 / (1.6\%) & \\
\hline
Hispanic/Latino(a) & 0 / (0\%) & 0 / (0\%) & 1 / (1.6\%) & \\
\hline
\_\_School Year\_\_ &  &  &  & 0.2\\
\hline
Freshman & 47 / (78\%) & 52 / (84\%) & 53 / (84\%) & \\
\hline
Sophmore & 8 / (13\%) & 6 / (9.7\%) & 8 / (13\%) & \\
\hline
Junior & 5 / (8.3\%) & 1 / (1.6\%) & 2 / (3.2\%) & \\
\hline
Senior & 0 / (0\%) & 3 / (4.8\%) & 0 / (0\%) & \\
\hline
\end{tabular}

\hypertarget{procedure}{%
\subsection{Procedure}\label{procedure}}

Participants were randomly assigned to one of three conditions representing different exposure to health benefits information. Our three conditions were an `active' intervention condition (n=60), a `passive' intervention condition (n=62), and our control condition (n=63). The two intervention conditions consisted of a packet of exercises adapted from the Choosing Healthplans All Together (CHAT) paradigm developed by Danis, Biddle \& Goold (2002). CHAT is a simulation exercise where participants construct their own HBP by allocating a limited set of resources to benefit types (e.g.~dental) and choosing scope of coverage (basic-to-high). The HBP as a whole is represented by a `game board' with several sections representing the different benefit types and with the scope of coverage represented by subdivisions in those sections. Each of these sections can be added to the HBP by paying a cost in markers representative of it's approximate relative cost in the US. For example, if a subject desired `basic' dental care (regular cleanings and examinations every 6 months, with minimal dental care), it would cost 2 markers. If the participant desired to upgrade to `medium' dental care (everything in `basic,' plus complete dental care including repairs and crowns), that would cost 4 additional markers, bringing the total cost to 6 markers. Participants have a total of 47 markers to use to design their HBP. Trade-offs are enforced as complete coverage is not possible with the limited resources. The core of the exercise is a simplified version of choosing priorities for a health care system.

Our active intervention condition had participants creating their own HBP through the CHAT exercise, while our passive intervention condition had participants being given a completed CHAT exercise. The HBP in our passive intervention condition, consisted of the consensus choices for health insurance found by Danis et al.~(2002). Our control condition was similar to the active intervention condition, but mentions of health care are replaced with pizza topping choices instead. Trade-offs are still enforced due to limited resources. This results in an exercise of similar length and intensity that is intentionally uninformative; see Appendix {[}LETTER HERE{]} for Study 1 active intervention materials, and {[}LETTER HERE{]} for passive intervention materials. Study 1 used a 2 (pre-post) x 3 (condition) mixed-subjects design, where condition was a between-subjects factor and participant were assigned to one of the three conditions. Time was a within-subjects factor with the primary outcome, support for UHC, measured before and after participants completed the control or one of the two intervention conditions.

\hypertarget{measures}{%
\subsection{Measures}\label{measures}}

The primary outcome measure was the support for UHC scale, adapted from Shen \& Labouff (2013), measured both pre and post-test. The scale was comprised of 4 items measuring support for UHC, which were averaged after reverse scoring the third item (e.g.~``Access to medical care and insurance is a basic, inherent right of man'') .Each item was measured on a 7 point Likert scale from 1 (strongly disagree) to 7 (strongly agree); see Table {[}LETTER HERE{]} for item wording.

Participants also responded to several items about their experience with health. Participants were asked whether they paid for their own health insurance and if they had ever been uninsured. Participants in the active intervention condition were also asked if they would be happy having the plan they built as their own health insurance. Each of these three items was measured as a `yes' or `no' response. Additionally, there was a free-response question asking about the subjects thoughts about the exercise they just completed. Finally, we also measured demographic information, including gender identity, age, race/ethnicity, and year in school.

\hypertarget{power-and-statistical-analyses}{%
\subsection{Power and Statistical Analyses}\label{power-and-statistical-analyses}}

We planned to recruit 180 participants. Sample size was determined a-priori using G-power with the following parameters: greater than 90\% power to determine a significant large-sized effect (Cohen's f = 0.10) at an alpha level of .05, for a linear multiple regression. Support for UHC outcome was treated as a continuous variable. We examined the effects of experimental condition (active intervention, passive intervention, and control ) and time (pre vs.~post) on our outcome variable by testing multi-level models with random and fixed intercepts. The linear mixed model we constructed had condition, time, and the condition x time interaction as our fixed effects. A random intercept for each of the subjects was included to account for within-subject correlation in scores. We examined the main effect and the 2-way interaction between our two predictors. Additionally, we also All tests were conducted in R and were considered statistically significant when P \textless{} .05.

Additionally, we fitted Bayesian linear multivariate multilevel models to our support for UHC outcome variable as a function of dummy-coded factors `condition' (reference level `control'), and `time' (reference level `pre') as well as the `condition x time' two way interaction using the Stan modeling language and the R package brms. Condition, time, and their interaction were our fixed effects, with a random intercept for subjects as our random effect. Our priors were a normal distribution with a mean of 0 and a standard deviation of 2.5 for the mean of our reference levels for our three fixed effects. We used the brms package's default priors for standard deviations of our random effects (Student's t-distribution with nu = 3, mu = 0 and sigma = 20), as well as for correlation coefficients in interaction models.

\hypertarget{study-1-hypothesis}{%
\subsection{Study 1 Hypothesis:}\label{study-1-hypothesis}}

Hypothesis 1 -- The experimental groups will differ in support for UHC.

\begin{quote}
H1a: Participants in the two intervention conditions will have greater increases in support for UHC compared to those in the control condition. We believe this to be the case due to HBPs directly addressing several common sources for opposition to UHC.
\end{quote}

\begin{quote}
H1b: Participants in the active intervention condition will have greater increases in support for UHC than participants in the passive intervention condition. We believe this to be the case as previous research indicates that complex, subject specific, numerical information is more easily learned through active engagement with the material.
\end{quote}

\hypertarget{results}{%
\section{Results}\label{results}}

\#note should we make sure to actually directly ref these results?

Descriptive statistics are summarized in {[}Table here{]}. Our hypothesis was tested using a linear mixed model fitted to our support for UHC outcome measure. Cronbach's alpha for the items in this measure was 0.85. In opposition to H1a and H1b, we observed no statistically significant effect in our planned comparison of our active intervention condition t(198.5) = 1.22 ,p = .224, or our passive intervention condition t(198.5) = 1.04 ,p = .299. Additionally, we observed no statistically significant effect in our planned comparison of time t(181) = 1.00 ,p = .317. Finally, we also saw no significant interaction between time and the active condition t(181) = 1.14 ,p = .258, or the passive condition t(181) = 1.67 ,p = .0963.

\begin{table}
\begin{center}
\begin{tabular}{l c}
\hline
 & Multi Level Model - Intercept Varies by Subject \\
\hline
Control                                                       & $4.79^{*}$       \\
                                                              & $ [ 4.51; 5.07]$ \\
Active Intervention                                           & $0.25$           \\
                                                              & $ [-0.15; 0.64]$ \\
Passive Intervention                                          & $0.21$           \\
                                                              & $ [-0.18; 0.60]$ \\
Post-Measurement Effect                                       & $0.06$           \\
                                                              & $ [-0.06; 0.18]$ \\
Interaction between Active Intervention and Post-Measurement  & $0.10$           \\
                                                              & $ [-0.07; 0.27]$ \\
Interaction between Passive Intervention and Post-Measurement & $0.14$           \\
                                                              & $ [-0.02; 0.31]$ \\
\hline
AIC                                                           & $830.98$         \\
BIC                                                           & $862.25$         \\
Log Likelihood                                                & $-407.49$        \\
Num. obs.                                                     & $368$            \\
Num. groups: SUBJECT                                          & $184$            \\
Var: SUBJECT (Intercept)                                      & $1.12$           \\
Var: Residual                                                 & $0.11$           \\
\hline
\multicolumn{2}{l}{\scriptsize{$^*$ 0 outside the confidence interval.}}
\end{tabular}
\caption{Frequentist Table of Intervention on UHC Support}
\label{tab:study_1_tabs}
\end{center}
\end{table}

For our Bayesian estimation, we had four sampling chains, each with 2000 iterations and 1000 warmup repetitions. This yielded 4000 estimated samples at convergence. Participants in our uninformative control condition had no significant change in support for UHC post intervention (℮ = 4.78, CI =4.49, 5.07) than pre intervention (℮ = 4.84, CI =4.55, 5.13). Participants in our `active' experimental condition had no difference in support for UHC post intervention (℮ = 5.03, CI =4.74, 5.32) than pre intervention (℮ =5.19, CI =4.90, 5.48). Participants in our `passive' experimental condition had no difference in support for UHC post intervention (℮ =4.99, CI =4.70, 1.34) versus pre intervention (℮ = 5.21 , CI =4.92, 5.50). In support of H1a, participants in both intervention conditions had greater support for UHC compared to the control. However, in opposition to H1b, participants in our active intervention condition did not have a greater increase in support for UHC compared to our passive condition.

\begin{figure}
\centering
\includegraphics{primary_files/figure-latex/study1_plot1-1.pdf}
\caption{(\#fig:study1\_plot1)Boxplot showing effect of interventions on support for UHC}
\end{figure}

\hypertarget{qualitative-results}{%
\subsection{Qualitative results}\label{qualitative-results}}

Analyzing our free-response question, we found several positive and negative trends in our findings. Some participants felt that the pencil and paper exercise was unnecessarily complex, and that the process of completing it was not self-explanitory. Several occasions occurred wherein the participant asked the administrator how to complete the exercise, after having read through the instructions. In total, 10\% of participants expressed some form of confusion in their free response segment, e.g.~
\textgreater{} ``It was a little confusing if you aren't very familiar with insurance and health care.''

\begin{quote}
``The way that plan was laid out with the pegs was slightly confusing and I think that it might provide more accurate answers if it were formatted more clearly. Otherwise, I thought that the different levels of care were described well and gave a good picture of what would be provided.''
\end{quote}

However, a larger proportion of our participants also found the exercise particularly engaging, interesting, fun, and helpful. In total, 32\% of participants expressed some form of positive engagement with the intervention conditions in their free response segment, e.g.

\begin{quote}
``I think that is was a good exercise to see what kind of benefits you would want and think about what benefits other people should have also''
\end{quote}

\begin{quote}
``It was fun trying to make those decisions. I ended up not bubbling any of the `long term'\,''retired person" options because I just don't care enough for those benefits. Maybe it's because I'm not in that situation yet.''
\end{quote}

We did not predict a priori that our intervention conditions would increase confusion. Nonetheless, an even larger contingent of our participants expressed positive feelings regarding the exercise. Given that the purpose of the interventions were to increase engagement with the often-times boring information necessary to explain UHC, our qualitative data indicates a positive outcome.

\hypertarget{discussion}{%
\section{Discussion}\label{discussion}}

The purpose of Study 1 was to determine if an HBP based intervention would lead to increased support for UHC (H1a). Our secondary aim was to determine if our active intervention condition would have a greater increase in support for UHC than participants in our passive intervention condition (H1b). The results of our linear-mixed model did not support either of our two hypotheses, as the HBP based intervention did not lead to increased support for UHC, and our active and passive intervention conditions had indistinguishable impact. Our Bayesian estimation provided support of our first hypothesis (H1a), indicating that participants in both our intervention conditions had greater support for UHC as compared to our control. However, our Bayesian estimation did not support our second hypothesis (H1b); Participants in our active intervention condition did not have a greater increase in support for UHC compared to our passive condition.

One plausible explanation regarding our conflicting results is the confusion regarding the experimental procedure and materials for participants in our active intervention condition. Primarily, our qualitative free-response section indicated that a significant portion of participants in our active intervention condition did not fully understand the instructions necessary. Considering the complexity and numerical engagement necessary to participate, a lack of understanding could plausibly blunt the impact of the intervention. Furthermore, if participants in our active intervention exhibited confusion, while participants in our passive intervention did not, as they had much simpler instructions, that would be a confounding variable when attempting to determine if active instruction is superior for communicating on UHC as compared to passive instruction.

A significant portion of our participants had expressed difficulties with the experimental protocol and materials. This adds avoidable stochasticity to our data and hinders reproducibility. Therefore, the single largest priority in moving from Study 1 to Study 2 was improving the experimental design and materials. Executing a pseudo-replication of the study, with a protocol that is designed to reduce confusion would significantly reduce potential confounding variables. Additionally, our control condition in Study 1 was an uninformative control, which is not necessarily a realistic comparison point regarding commonly available information on UHC. Thus, the second priority was to alter our control condition to reflect the messaging more accurately on UHC that is already available to improve external validity.

\hypertarget{method-1}{%
\section{Method}\label{method-1}}

\hypertarget{participants-1}{%
\subsection{Participants}\label{participants-1}}

Our participants were 189 students enrolled in a Psychology course at a Midwestern University. Our participants were primarily white (76\%), female (68\%), and freshmen (80\%); further demographic information can be found in {[}Table here{]}. Participants received course credit for participation in the study.

\begin{tabular}{l|l|l|l|l}
\hline
**Characteristic** & **Active Intervention**, N = 60 & **No Intervention**, N = 62 & **Passive Intervention**, N = 63 & **p-value**\\
\hline
\_\_Age\_\_ &  &  &  & 0.10\\
\hline
18 & 34 / (57\%) & 46 / (74\%) & 44 / (70\%) & \\
\hline
19 & 19 / (32\%) & 10 / (16\%) & 16 / (25\%) & \\
\hline
20 & 7 / (12\%) & 3 / (4.8\%) & 2 / (3.2\%) & \\
\hline
21 & 0 / (0\%) & 2 / (3.2\%) & 1 / (1.6\%) & \\
\hline
23 & 0 / (0\%) & 1 / (1.6\%) & 0 / (0\%) & \\
\hline
\_\_Sex\_\_ &  &  &  & 0.7\\
\hline
Female & 36 / (60\%) & 41 / (66\%) & 42 / (67\%) & \\
\hline
Male & 24 / (40\%) & 21 / (34\%) & 21 / (33\%) & \\
\hline
\_\_Race\_\_ &  &  &  & >0.9\\
\hline
Caucasian/White & 46 / (77\%) & 49 / (79\%) & 48 / (76\%) & \\
\hline
African-American/Black & 5 / (8.3\%) & 5 / (8.1\%) & 6 / (9.5\%) & \\
\hline
Asian/Pacific Islander & 3 / (5.0\%) & 3 / (4.8\%) & 3 / (4.8\%) & \\
\hline
Caucasian/White,Asian/Pacific Islander & 2 / (3.3\%) & 4 / (6.5\%) & 1 / (1.6\%) & \\
\hline
Caucasian/White,African-American/Black & 1 / (1.7\%) & 1 / (1.6\%) & 1 / (1.6\%) & \\
\hline
 & 2 / (3.3\%) & 0 / (0\%) & 0 / (0\%) & \\
\hline
Other & 1 / (1.7\%) & 0 / (0\%) & 1 / (1.6\%) & \\
\hline
African-American/Black,American Indian/Alaska Native & 0 / (0\%) & 0 / (0\%) & 1 / (1.6\%) & \\
\hline
African-American/Black,Hispanic/Latino(a) & 0 / (0\%) & 0 / (0\%) & 1 / (1.6\%) & \\
\hline
Hispanic/Latino(a) & 0 / (0\%) & 0 / (0\%) & 1 / (1.6\%) & \\
\hline
\_\_School Year\_\_ &  &  &  & 0.2\\
\hline
Freshman & 47 / (78\%) & 52 / (84\%) & 53 / (84\%) & \\
\hline
Sophmore & 8 / (13\%) & 6 / (9.7\%) & 8 / (13\%) & \\
\hline
Junior & 5 / (8.3\%) & 1 / (1.6\%) & 2 / (3.2\%) & \\
\hline
Senior & 0 / (0\%) & 3 / (4.8\%) & 0 / (0\%) & \\
\hline
\end{tabular}

\hypertarget{procedure-1}{%
\subsection{Procedure}\label{procedure-1}}

Participants were randomly assigned to one of three conditions representing different exposure to health benefits information. Our three conditions were an `active' intervention condition (n=60), a `passive' intervention condition (n=62), and our control condition (n=63). The two intervention conditions consisted of a packet of exercises adapted from the Choosing Healthplans All Together (CHAT) paradigm developed by Danis, Biddle \& Goold (2002). CHAT is a simulation exercise where participants construct their own HBP by allocating a limited set of resources to benefit types (e.g.~dental) and choosing scope of coverage (basic-to-high). The HBP as a whole is represented by a `game board' with several sections representing the different benefit types and with the scope of coverage represented by subdivisions in those sections. Each of these sections can be added to the HBP by paying a cost in markers representative of it's approximate relative cost in the US. For example, if a subject desired `basic' dental care (regular cleanings and examinations every 6 months, with minimal dental care), it would cost 2 markers. If the participant desired to upgrade to `medium' dental care (everything in `basic,' plus complete dental care including repairs and crowns), that would cost 4 additional markers, bringing the total cost to 6 markers. Participants have a total of 47 markers to use to design their HBP. Trade-offs are enforced as complete coverage is not possible with the limited resources. The core of the exercise is a simplified version of choosing priorities for a health care system.

Our active intervention condition had participants creating their own HBP through the CHAT exercise, while our passive intervention condition had participants being given a completed CHAT exercise. The HBP in our passive intervention condition, consisted of the consensus choices for health insurance found by Danis et al.~(2002). Our control condition was similar to the active intervention condition, but mentions of health care are replaced with pizza topping choices instead. Trade-offs are still enforced due to limited resources. This results in an exercise of similar length and intensity that is intentionally uninformative; see Appendix {[}LETTER HERE{]} for Study 1 active intervention materials, and {[}LETTER HERE{]} for passive intervention materials. Study 1 used a 2 (pre-post) x 3 (condition) mixed-subjects design, where condition was a between-subjects factor and participant were assigned to one of the three conditions. Time was a within-subjects factor with the primary outcome, support for UHC, measured before and after participants completed the control or one of the two intervention conditions.

\hypertarget{measures-1}{%
\subsection{Measures}\label{measures-1}}

The primary outcome measure was the support for UHC scale, adapted from Shen \& Labouff (2013), measured both pre and post-test. The scale was comprised of 4 items measuring support for UHC, which were averaged after reverse scoring the third item (e.g.~``Access to medical care and insurance is a basic, inherent right of man'') .Each item was measured on a 7 point Likert scale from 1 (strongly disagree) to 7 (strongly agree); see Table {[}LETTER HERE{]} for item wording.

Participants also responded to several items about their experience with health. Participants were asked whether they paid for their own health insurance and if they had ever been uninsured. Participants in the active intervention condition were also asked if they would be happy having the plan they built as their own health insurance. Each of these three items was measured as a `yes' or `no' response. Additionally, there was a free-response question asking about the subjects thoughts about the exercise they just completed. Finally, we also measured demographic information, including gender identity, age, race/ethnicity, and year in school.

\hypertarget{power-and-statistical-analyses-1}{%
\subsection{Power and Statistical Analyses}\label{power-and-statistical-analyses-1}}

We planned to recruit 180 participants. Sample size was determined a-priori using G-power with the following parameters: greater than 90\% power to determine a significant large-sized effect (Cohen's f = 0.10) at an alpha level of .05, for a linear multiple regression. Support for UHC outcome was treated as a continuous variable. We examined the effects of experimental condition (active intervention, passive intervention, and control ) and time (pre vs.~post) on our outcome variable by testing multi-level models with random and fixed intercepts. The linear mixed model we constructed had condition, time, and the condition x time interaction as our fixed effects. A random intercept for each of the subjects was included to account for within-subject correlation in scores. We examined the main effect and the 2-way interaction between our two predictors. Additionally, we also All tests were conducted in R and were considered statistically significant when P \textless{} .05.

Additionally, we fitted Bayesian linear multivariate multilevel models to our support for UHC outcome variable as a function of dummy-coded factors `condition' (reference level `control'), and `time' (reference level `pre') as well as the `condition x time' two way interaction using the Stan modeling language and the R package brms. Condition, time, and their interaction were our fixed effects, with a random intercept for subjects as our random effect. Our priors were a normal distribution with a mean of 0 and a standard deviation of 2.5 for the mean of our reference levels for our three fixed effects. We used the brms package's default priors for standard deviations of our random effects (Student's t-distribution with nu = 3, mu = 0 and sigma = 20), as well as for correlation coefficients in interaction models.

\hypertarget{study-1-hypothesis-1}{%
\subsection{Study 1 Hypothesis:}\label{study-1-hypothesis-1}}

Hypothesis 1 -- The experimental groups will differ in support for UHC.

\begin{quote}
H1a: Participants in the two intervention conditions will have greater increases in support for UHC compared to those in the control condition. We believe this to be the case due to HBPs directly addressing several common sources for opposition to UHC.
\end{quote}

\begin{quote}
H1b: Participants in the active intervention condition will have greater increases in support for UHC than participants in the passive intervention condition. We believe this to be the case as previous research indicates that complex, subject specific, numerical information is more easily learned through active engagement with the material.
\end{quote}

\hypertarget{discussion-1}{%
\section{Discussion}\label{discussion-1}}

The purpose of Study 1 was to determine if an HBP based intervention would lead to increased support for UHC (H1a). Our secondary aim was to determine if our active intervention condition would have a greater increase in support for UHC than participants in our passive intervention condition (H1b). The results of our linear-mixed model did not support either of our two hypotheses, as the HBP based intervention did not lead to increased support for UHC, and our active and passive intervention conditions had indistinguishable impact. Our Bayesian estimation provided support of our first hypothesis (H1a), indicating that participants in both our intervention conditions had greater support for UHC as compared to our control. However, our Bayesian estimation did not support our second hypothesis (H1b); Participants in our active intervention condition did not have a greater increase in support for UHC compared to our passive condition.

One plausible explanation regarding our conflicting results is the confusion regarding the experimental procedure and materials for participants in our active intervention condition. Primarily, our qualitative free-response section indicated that a significant portion of participants in our active intervention condition did not fully understand the instructions necessary. Considering the complexity and numerical engagement necessary to participate, a lack of understanding could plausibly blunt the impact of the intervention. Furthermore, if participants in our active intervention exhibited confusion, while participants in our passive intervention did not, as they had much simpler instructions, that would be a confounding variable when attempting to determine if active instruction is superior for communicating on UHC as compared to passive instruction.

A significant portion of our participants had expressed difficulties with the experimental protocol and materials. This adds avoidable stochasticity to our data and hinders reproducibility. Therefore, the single largest priority in moving from Study 1 to Study 2 was improving the experimental design and materials. Executing a pseudo-replication of the study, with a protocol that is designed to reduce confusion would significantly reduce potential confounding variables. Additionally, our control condition in Study 1 was an uninformative control, which is not necessarily a realistic comparison point regarding commonly available information on UHC. Thus, the second priority was to alter our control condition to reflect the messaging more accurately on UHC that is already available to improve external validity.

\texttt{\#\{r\ child\ =\ "appendix\_a.Rmd"\}}

\hypertarget{appendix-b}{%
\section{Appendix B}\label{appendix-b}}

\hypertarget{experimental-materials-for-study-2}{%
\subsubsection{Experimental Materials for Study 2}\label{experimental-materials-for-study-2}}

\textbackslash begin\{figure\}
\includegraphics[width=1\linewidth]{appendix_b_1}

\caption{Screenshot of Web Exercise for Intervention Condition}

(\#fig:appendix\_b1)
\textbackslash end\{figure\}

\textbackslash begin\{figure\}
\includegraphics[width=1\linewidth]{appendix_b_2}

\caption{Screenshot showing legend and web exercise for intervention condition}

(\#fig:appendix\_b2)
\textbackslash end\{figure\}

\textbackslash begin\{figure\}
\includegraphics[width=1\linewidth]{appendix_b_3}

\caption{1st Infographic for control condition}

(\#fig:appendix\_b3)
\textbackslash end\{figure\}

\textbackslash begin\{figure\}
\includegraphics[width=1\linewidth]{appendix_b_4}

\caption{2nd Infographic for control condition}

(\#fig:appendix\_b4)
\textbackslash end\{figure\}

\textbackslash begin\{figure\}
\includegraphics[width=1\linewidth]{appendix_b_5}

\caption{3rd Infographic for control condition}

(\#fig:appendix\_b5)
\textbackslash end\{figure\}

\textbackslash begin\{figure\}
\includegraphics[width=1\linewidth]{appendix_b_6}

\caption{Support for UHC Measure - Scale and Item wording}

(\#fig:appendix\_b6)
\textbackslash end\{figure\}

\textbackslash begin\{figure\}
\includegraphics[width=1\linewidth]{appendix_b_7}

\caption{Percieved Equity Measure - Scale and Item wording}

(\#fig:appendix\_b7)
\textbackslash end\{figure\}

\textbackslash begin\{figure\}
\includegraphics[width=1\linewidth]{appendix_b_8}

\caption{Percieved Understanding Measure - Scale and Item wording}

(\#fig:appendix\_b8)
\textbackslash end\{figure\}

\textbackslash begin\{figure\}
\includegraphics[width=1\linewidth]{appendix_b_9}

\caption{Subjective Numeracy Scale Part 1 - Scale and Item wording}

(\#fig:appendix\_b9)
\textbackslash end\{figure\}

\textbackslash begin\{figure\}
\includegraphics[width=1\linewidth]{appendix_b_10}

\caption{Subjective Numeracy Scale Part 2 - Scale and Item wording}

(\#fig:appendix\_b10)
\textbackslash end\{figure\}

\textbackslash begin\{figure\}
\includegraphics[width=1\linewidth]{appendix_b_11}

\caption{Rasch Numeracy Scale (Objective Numeracy) - Item wording}

(\#fig:appendix\_b11)
\textbackslash end\{figure\}


\end{document}
